\documentclass[../../book.tex]{subfiles}

\begin{document}

\subsection{Rejection Sampling and Lyubashevsky Signatures}\label{section:lyubashevsky}

With SIS and LWE at hand, we can build lattice-based analogues of the discrete-logarithm protocols from the previous chapters.
In this section, we translate the Schnorr signature into the lattice setting.
The translation fails in an instructive way, and the fix, called \textbf{rejection sampling}, is the central technique of lattice-based signatures and zero-knowledge proofs.
The resulting scheme is due to Lyubashevsky~\cite{EC:Lyubashevsky12}, and the same idea underlies ML-DSA (Dilithium).

For simplicity, we work over $\mathbb{Z}_q$; the ring and module versions are obtained by replacing $\mathbb{Z}_q$ with $\mathcal{R}_q$.
Besides $S_{\eta}$, we use the challenge set
\begin{equation*}
    B_{\kappa} := \{\mathbf{c} \in \{-1,0,1\}^k: \mathbf{c} \; \text{has exactly} \; \kappa \; \text{non-zero entries}\}, \quad |B_{\kappa}| = 2^{\kappa}\binom{k}{\kappa},
\end{equation*}
where $\kappa$ is chosen so that $|B_{\kappa}| \geq 2^{\lambda}$.

\subsubsection{From Schnorr to Lattices}

Recall the Schnorr signature over a group $\mathbb{G} = \langle g \rangle$ of order $q$.
The secret key is $\alpha \in \mathbb{Z}_q$ and the public key is $h = g^{\alpha}$.
To sign $\mu$, the signer samples $r \leftarrowS \mathbb{Z}_q$, computes the challenge $c \gets \mathsf{H}(\mu, g^r)$, and outputs $(z := r + \alpha c, c)$.
The verifier recomputes $g^r = g^z h^{-c}$ and checks that $c = \mathsf{H}(\mu, g^zh^{-c})$.
Two properties make this scheme secure: $\alpha \mapsto g^{\alpha}$ is one-way, and $z$ is uniform in $\mathbb{Z}_q$ and thus reveals nothing about $\alpha$.

\paragraph{Attempt \#1: replace exponentiation with a matrix}
The lattice analogue of $\alpha \mapsto g^{\alpha}$ is the linear map $\mathbf{S} \mapsto \mathbf{A}\mathbf{S}$ for a public random matrix $\mathbf{A} \in \mathbb{Z}_q^{n \times m}$.
The secret key becomes $\mathbf{S} \in \mathbb{Z}_q^{m \times k}$, the public key becomes $\mathbf{T} := \mathbf{A}\mathbf{S}$, and the signature becomes $(\mathbf{z} := \mathbf{r} + \mathbf{S}\mathbf{c}, \mathbf{c})$ for $\mathbf{c} \gets \mathsf{H}(\mu, \mathbf{A}\mathbf{r})$.
The verifier checks $\mathbf{c} = \mathsf{H}(\mu, \mathbf{A}\mathbf{z} - \mathbf{T}\mathbf{c})$.

This scheme is completely insecure.
Since $m > n$, anyone can solve the linear system $\mathbf{A}\mathbf{S}' = \mathbf{T}$ by Gaussian elimination and obtain a (long) matrix $\mathbf{S}'$ that works as a secret key.
The map $\mathbf{S} \mapsto \mathbf{A}\mathbf{S}$ is one-way only on \emph{short} inputs, by ISIS.

\paragraph{Attempt \#2: make everything short}
We therefore sample $\mathbf{S} \leftarrowS S_{\eta}^{m \times k}$ and require the verifier to check that $\mathbf{z}$ is short, so that a forger cannot use a long $\mathbf{S}'$.
For an honest $\mathbf{z} = \mathbf{r} + \mathbf{S}\mathbf{c}$ to pass this check, both terms must be short.
The term $\mathbf{S}\mathbf{c}$ is short by construction: each of its entries is a sum of $\kappa$ entries of $\mathbf{S}$ (up to signs), so $\|\mathbf{S}\mathbf{c}\|_{\infty} \leq \eta\kappa$.
The masking vector $\mathbf{r}$ is then sampled from some short distribution, and all computations of $\mathbf{z}$ are done over $\mathbb{Z}$, without reduction modulo $q$.

This attempt is still insecure, but for a subtler reason: \emph{$\mathbf{z}$ leaks the secret key}.
A short $\mathbf{r}$ cannot hide the shift $\mathbf{S}\mathbf{c}$ completely.
For example, if $\mathbf{r}$ is centered at zero, then $\mathbb{E}[\mathbf{z}] = \mathbf{S}\mathbf{c}$, and since $\mathbf{c}$ is public, an observer who collects enough signatures recovers $\mathbf{S}$ by linear regression.
In Schnorr signatures this problem never arises, because $r$ is uniform over the whole group $\mathbb{Z}_q$ and $r + \alpha c$ is uniform for every $\alpha$.

\paragraph{Attempt \#3: rejection sampling}
The fix is to \emph{abort} and restart the signing procedure whenever the candidate $\mathbf{z}$ reveals too much about $\mathbf{S}$.
The abort probability is chosen so that the distribution of every \emph{output} $\mathbf{z}$ is a fixed distribution independent of $\mathbf{S}$.
Since signing is the Fiat--Shamir transform (\Cref{subsection:fiat-shamir}) of an identification protocol, this approach is called \textbf{Fiat--Shamir with aborts}.
The tool that computes the right abort probability is the following lemma.

\begin{lemma}[Rejection Sampling]\label{lemma:rejection-sampling}
    Let $f$ and $g$ be probability distributions over a set $V$ such that $f(z) \leq M g(z)$ for all $z \in V$ and some constant $M \geq 1$.
    Consider the following procedure:
    \begin{algoen}
        \item sample $z \leftarrowS g$;
        \item with probability $\frac{f(z)}{Mg(z)}$, \textbf{output} $z$; otherwise, go back to step 1.
    \end{algoen}
    The output of this procedure is distributed exactly according to $f$, and the expected number of iterations is $M$.
\end{lemma}

\begin{proof}
    In a single iteration, the procedure samples a particular $z$ and accepts it with probability $g(z) \cdot \frac{f(z)}{Mg(z)} = \frac{f(z)}{M}$.
    Summing over all $z$, one iteration succeeds with probability $\frac{1}{M}$, so the expected number of iterations is $M$.
    Conditioned on success, $z$ is output with probability $\frac{f(z)/M}{1/M} = f(z)$.
\end{proof}

\begin{intuition}
    Think of the ``graph'' of $Mg$ as a tent that covers the graph of $f$ everywhere.
    Sampling from $g$ and accepting with probability $f(z)/Mg(z)$ is equivalent to throwing a uniformly random point under the tent and keeping it only if it also lies under the graph of $f$.
    The kept points are uniformly distributed under the graph of $f$, and so their $z$-coordinates are distributed according to $f$.
    The ratio of the two areas is $M$, which is the expected number of throws.
\end{intuition}

In the signature scheme, $g$ is the distribution of the candidate $\mathbf{z} = \mathbf{r} + \mathbf{S}\mathbf{c}$ (which depends on $\mathbf{S}\mathbf{c}$), and $f$ is a fixed target distribution that does not depend on $\mathbf{S}$.
The lemma requires $f \leq Mg$ for a small constant $M$ and \emph{every} possible shift $\mathbf{S}\mathbf{c}$, and this determines how to choose the distribution of $\mathbf{r}$.
We now present two standard choices.

\subsubsection{Rejection Sampling with Uniform Distributions}

The simplest choice, used in ML-DSA, samples $\mathbf{r}$ uniformly from a box $[-\gamma,\gamma]^m$ that is much wider than the shift.
Fix a shift $\mathbf{v} := \mathbf{S}\mathbf{c}$ with $\|\mathbf{v}\|_{\infty} \leq \beta := \eta\kappa$ and consider a single coordinate $z_i = r_i + v_i$.
This coordinate is uniform over the interval $[-\gamma+v_i, \gamma+v_i]$, which depends on $v_i$ but always contains the smaller interval $[-(\gamma-\beta), \gamma-\beta]$.

Accept $\mathbf{z}$ if and only if $\|\mathbf{z}\|_{\infty} \leq \gamma - \beta$.
Conditioned on acceptance, $\mathbf{z}$ is uniform over the box $[-(\gamma-\beta),\gamma-\beta]^m$, which does not depend on $\mathbf{v}$.
This is precisely \Cref{lemma:rejection-sampling} with $f$ uniform over the small box, $g$ uniform over the shifted large box, and $M = \left(\frac{2\gamma+1}{2(\gamma-\beta)+1}\right)^m$.
The acceptance probability is
\begin{equation*}
    \frac{1}{M} = \left(\frac{2(\gamma-\beta)+1}{2\gamma+1}\right)^m \approx \left(1 - \frac{\beta}{\gamma}\right)^m \approx e^{-m\beta/\gamma},
\end{equation*}
so choosing $\gamma \approx m\beta$ gives $M \approx e$, \textit{i.e.}, fewer than three attempts per signature on average.

\begin{exercise}
    Let $m = 1$, $\beta = 1$, and $\gamma = 3$.
    List the distribution of $z = r + v$ for $v = -1$, $v = 0$, and $v = 1$, and check that after rejecting $|z| > 2$ all three distributions coincide.
\end{exercise}

\subsubsection{Rejection Sampling with Discrete Gaussians}

The uniform box forces $\gamma$ to grow linearly with $m$.
Lyubashevsky's original scheme instead uses the discrete Gaussian distribution, which leads to a width that grows only with $\sqrt{m}$ and therefore to shorter signatures.

\begin{definition}[Discrete Gaussian]
    For $\sigma > 0$ and $\mathbf{v} \in \mathbb{Z}^m$, let $\rho_{\mathbf{v},\sigma}(\mathbf{x}) := \exp\left(-\frac{\|\mathbf{x}-\mathbf{v}\|_2^2}{2\sigma^2}\right)$ for $\mathbf{x} \in \mathbb{R}^m$.
    The \textbf{discrete Gaussian distribution} over $\mathbb{Z}^m$ with center $\mathbf{v}$ and width $\sigma$ assigns to each $\mathbf{x} \in \mathbb{Z}^m$ the probability
    \begin{equation*}
        D^m_{\mathbf{v},\sigma}(\mathbf{x}) := \frac{\rho_{\mathbf{v},\sigma}(\mathbf{x})}{\sum_{\mathbf{y} \in \mathbb{Z}^m}\rho_{\mathbf{0},\sigma}(\mathbf{y})}.
    \end{equation*}
    We omit $\mathbf{v}$ when $\mathbf{v} = \mathbf{0}$ and write $D^m_{\sigma}$.
\end{definition}

The normalization is the same for every integer center $\mathbf{v}$, since shifting by $\mathbf{v}$ permutes $\mathbb{Z}^m$.
Informally, $D^m_{\mathbf{v},\sigma}$ is the usual bell-shaped Gaussian centered at $\mathbf{v}$, restricted to integer points\footnote{Some works, including parts of the lattice literature, write the Gaussian as $\exp(-\pi\|\mathbf{x}\|_2^2/s^2)$. The two conventions are related by $s = \sqrt{2\pi}\sigma$.}.
We need two facts about it.
First, samples are short: for $\mathbf{z} \leftarrowS D^m_{\sigma}$ and any $\tau > 1$, one has $\Pr[\|\mathbf{z}\|_2 > \tau\sigma\sqrt{m}] \leq \tau^m e^{m(1-\tau^2)/2}$~\cite[Lemma 4.4]{EC:Lyubashevsky12}, which is negligible already for $\tau = 2$.
Second, one-dimensional projections are also Gaussian: for any fixed $\mathbf{v}$ and $t > 0$, one has $\Pr[|\langle \mathbf{z},\mathbf{v} \rangle| > t] \leq 2\exp\left(-\frac{t^2}{2\|\mathbf{v}\|_2^2\sigma^2}\right)$.

Now let $\mathbf{r} \leftarrowS D^m_{\sigma}$, so that the candidate $\mathbf{z} = \mathbf{r} + \mathbf{v}$ follows the shifted distribution $g := D^m_{\mathbf{v},\sigma}$.
The target is the centered distribution $f := D^m_{\sigma}$.
The ratio between them is
\begin{equation*}
    \frac{D^m_{\sigma}(\mathbf{z})}{D^m_{\mathbf{v},\sigma}(\mathbf{z})} = \exp\left(\frac{\|\mathbf{z}-\mathbf{v}\|_2^2 - \|\mathbf{z}\|_2^2}{2\sigma^2}\right) = \exp\left(\frac{\|\mathbf{v}\|_2^2 - 2\langle \mathbf{z},\mathbf{v} \rangle}{2\sigma^2}\right).
\end{equation*}
This ratio is \emph{not} bounded for all $\mathbf{z}$, since $\langle \mathbf{z},\mathbf{v} \rangle$ can be arbitrarily negative.
However, by the second fact with $t = 12\sigma\|\mathbf{v}\|_2$, the event $|\langle \mathbf{z},\mathbf{v} \rangle| \leq 12\sigma\|\mathbf{v}\|_2$ holds except with probability $2e^{-72} < 2^{-100}$.
On this event, writing $\sigma = \alpha\|\mathbf{v}\|_2$, we get
\begin{equation*}
    \frac{D^m_{\sigma}(\mathbf{z})}{D^m_{\mathbf{v},\sigma}(\mathbf{z})} \leq \exp\left(\frac{12\|\mathbf{v}\|_2}{\sigma} + \frac{\|\mathbf{v}\|_2^2}{2\sigma^2}\right) = \exp\left(\frac{12}{\alpha} + \frac{1}{2\alpha^2}\right) =: M.
\end{equation*}
The shift satisfies $\|\mathbf{S}\mathbf{c}\|_2 \leq \sqrt{m}\|\mathbf{S}\mathbf{c}\|_{\infty} \leq \eta\kappa\sqrt{m} =: T$.
Choosing $\sigma = 12T$ makes $\alpha \geq 12$ for every possible shift, which gives $M \leq e^{1+1/288} \approx 2.72$.

Therefore, the condition $f \leq Mg$ of \Cref{lemma:rejection-sampling} holds except with probability below $2^{-100}$.
A slight generalization of the lemma~\cite[Lemma 4.7]{EC:Lyubashevsky12} shows that this failure event only adds a negligible statistical distance between the output and $D^m_{\sigma}$.

\subsubsection{The Lyubashevsky Signature Scheme}

We are now ready to state the scheme.
The matrix $\mathbf{A} \leftarrowS \mathbb{Z}_q^{n \times m}$ is a public parameter, $\mathsf{H}: \{0,1\}^* \to B_{\kappa}$ is a hash function modelled as a random oracle, and $\sigma = 12\eta\kappa\sqrt{m}$ and $M = e^{1+1/288}$ are as computed above.

\begin{tcolorbox}[title=Lyubashevsky Signature Scheme,
    breakable,
    colback=blue!5!white,
    colframe=blue!75!black,
    colbacktitle=blue!25!white,
    coltitle=blue!20!black,
    fonttitle=\bfseries,
    boxrule=1.25pt,
    subtitle style={boxrule=0pt,
    colback=blue!20!white,
    colupper=blue!75!gray} ]

    \tcbsubtitle{$\mathsf{KeyGen}(1^{\lambda})$}
    \begin{itemize}[label=\ding{51}]
        \item Sample $\mathbf{S} \leftarrowS S_{\eta}^{m \times k}$ and compute $\mathbf{T} \gets \mathbf{A}\mathbf{S} \in \mathbb{Z}_q^{n \times k}$.
        \item \textbf{Output} $\mathsf{pk} \gets \mathbf{T}$ and $\mathsf{sk} \gets \mathbf{S}$.
    \end{itemize}
    \tcbsubtitle{$\mathsf{Sign}(\mathsf{sk}, \mu)$}
    \begin{itemize}[label=\ding{51}]
        \item Sample the mask $\mathbf{r} \leftarrowS D^m_{\sigma}$.
        \item Compute the challenge $\mathbf{c} \gets \mathsf{H}(\mu, \mathbf{A}\mathbf{r} \bmod q)$.
        \item Compute $\mathbf{z} \gets \mathbf{r} + \mathbf{S}\mathbf{c}$ over $\mathbb{Z}$ \commentgray{no reduction modulo $q$!}
        \item With probability $\min\left\{1, \frac{D^m_{\sigma}(\mathbf{z})}{M \cdot D^m_{\mathbf{S}\mathbf{c},\sigma}(\mathbf{z})}\right\}$, \textbf{output} $(\mathbf{z},\mathbf{c})$; otherwise, restart from the first step.
    \end{itemize}
    \tcbsubtitle{$\mathsf{Verify}(\mathsf{pk}, \mu, (\mathbf{z},\mathbf{c}))$}
    \begin{itemize}[label=\ding{51}]
        \item Check that $\|\mathbf{z}\|_2 \leq 2\sigma\sqrt{m}$.
        \item Check that $\mathbf{c} = \mathsf{H}(\mu, \mathbf{A}\mathbf{z} - \mathbf{T}\mathbf{c} \bmod q)$.
    \end{itemize}
\end{tcolorbox}

\textbf{Correctness.}
For an honest signature, $\mathbf{A}\mathbf{z} - \mathbf{T}\mathbf{c} = \mathbf{A}\mathbf{r} + \mathbf{A}\mathbf{S}\mathbf{c} - \mathbf{A}\mathbf{S}\mathbf{c} = \mathbf{A}\mathbf{r}$, so the hash check passes.
The output $\mathbf{z}$ is distributed (up to negligible statistical distance) as $D^m_{\sigma}$, so the norm check passes except with probability $2^m e^{-3m/2} = e^{-(3/2-\ln 2)m}$ by the first Gaussian fact.
The signer needs $M \approx 2.72$ attempts on average.

\textbf{Security (sketch).}
The proof follows the Schnorr proof and has two steps.
\begin{enumerate}
    \item[\textit{(i)}] \textbf{Simulation.} The reduction must answer signing queries without $\mathbf{S}$.
    Thanks to rejection sampling, it samples $\mathbf{c} \leftarrowS B_{\kappa}$ and $\mathbf{z} \leftarrowS D^m_{\sigma}$ directly, outputs them with probability $1/M$, and programs the random oracle as $\mathsf{H}(\mu, \mathbf{A}\mathbf{z}-\mathbf{T}\mathbf{c}) := \mathbf{c}$.
    By the rejection sampling lemma, these simulated signatures are statistically indistinguishable from real ones.
    Without rejection sampling, the real $\mathbf{z}$ would depend on $\mathbf{S}$ and this simulation would fail.
    \item[\textit{(ii)}] \textbf{Extraction.} By the forking lemma, a successful forger can be rewound to produce two valid signatures $(\mathbf{z},\mathbf{c})$ and $(\mathbf{z}',\mathbf{c}')$ with $\mathbf{c} \neq \mathbf{c}'$ and the same value $\mathbf{A}\mathbf{z}-\mathbf{T}\mathbf{c} = \mathbf{A}\mathbf{z}'-\mathbf{T}\mathbf{c}'$.
    Substituting $\mathbf{T} = \mathbf{A}\mathbf{S}$ gives
    \begin{equation*}
        \mathbf{A}\left(\mathbf{z}-\mathbf{z}' + \mathbf{S}(\mathbf{c}'-\mathbf{c})\right) = \mathbf{0} \pmod q,
    \end{equation*}
    and the vector in parentheses has $\ell^2$ norm at most $4\sigma\sqrt{m} + 2\eta\kappa\sqrt{m}$.
    It is a solution to SIS in the $\ell^2$ norm, provided it is non-zero\footnote{Non-zeroness is the subtle part of the proof. The parameter $m$ is chosen large enough that many secret keys $\mathbf{S}'$ with $\mathbf{A}\mathbf{S}' = \mathbf{T}$ exist, and the forger cannot tell which one the reduction holds. Hence, with probability at least $1/2$, the vector is non-zero for the key the reduction actually uses.}.
\end{enumerate}

\begin{remark}[Relaxed extraction]\label{remark:relaxed-extraction}
    Notice what the extractor in step \textit{(ii)} actually obtains: not the secret $\mathbf{S}$, but the pair $(\bar{\mathbf{z}}, \bar{\mathbf{c}}) := (\mathbf{z}-\mathbf{z}', \mathbf{c}-\mathbf{c}')$ with $\mathbf{A}\bar{\mathbf{z}} = \mathbf{T}\bar{\mathbf{c}}$.
    The extracted vector $\bar{\mathbf{z}}$ is about $\sigma/\eta$ times longer than the honest secret, and it satisfies the relation only after multiplying $\mathbf{T}$ by the difference of challenges $\bar{\mathbf{c}}$.
    This gap between what the honest prover knows and what the extractor can guarantee is inherent to all lattice-based proofs of knowledge (it is often called \emph{soundness slack}), and we will meet it again when building lattice-based zero-knowledge proofs.
\end{remark}

\paragraph{Performance}
For about $100$ bits of security, Lyubashevsky~\cite{EC:Lyubashevsky12} proposes, among others, the parameters $n = 512$, $q = 2^{27}$, $m = 8786$, $k = 80$, $\eta = 1$, and $\kappa = 28$, which give $\sigma \approx 31495$.
The resulting signatures take about $19.9$~KB and the keys about $128$~KB each, which is far from practical.
The size comes mostly from the unstructured matrix and the large dimension $m$.
Moving to Module-SIS and Module-LWE over $\mathcal{R}_q$ (\Cref{section:structured-lattices}), using the uniform rejection sampling, and compressing the public key and the signature yield ML-DSA (Dilithium)~\cite{TCHES:DKLLSS18}, standardized in FIPS~204~\cite{FIPS204}, with signatures of about $2.4$~KB and public keys of about $1.3$~KB at its lowest security level.

\end{document}